\documentclass[11pt,a4paper]{article}
\usepackage[utf8]{inputenc}
\usepackage[french]{babel}
\usepackage[T1]{fontenc}
\usepackage{amsmath,amssymb,amsfonts}
\usepackage{geometry}
\geometry{hmargin=2.5cm,vmargin=3cm}
\usepackage{graphicx}
\usepackage{booktabs}
\usepackage{listings}
\usepackage{xcolor}
\usepackage{hyperref}
\usepackage{tikz}
\usetikzlibrary{shapes.geometric, arrows, positioning}

% Configuration de l'affichage du code
\definecolor{codegreen}{rgb}{0,0.6,0}
\definecolor{codegray}{rgb}{0.5,0.5,0.5}
\definecolor{codepurple}{rgb}{0.58,0,0.82}
\definecolor{backcolour}{rgb}{0.96,0.96,0.94}

\lstdefinestyle{mystyle}{
    backgroundcolor=\color{backcolour},   
    commentstyle=\color{codegreen},
    keywordstyle=\color{blue}\bfseries,
    numberstyle=\tiny\color{codegray},
    stringstyle=\color{codepurple},
    basicstyle=\ttfamily\footnotesize,
    breakatwhitespace=false,         
    breaklines=true,                 
    captionpos=b,                    
    keepspaces=true,                 
    numbers=left,                    
    numbersep=5pt,                  
    showspaces=false,                
    showstringspaces=false,
    showtabs=false,                  
    tabsize=2
}
\lstset{style=mystyle}

\title{\textbf{Conception et implémentation d'un HIDS basé sur eBPF sous Linux \\ \large Monographie Universitaire et Technique du Projet Block-X}}
\author{\textbf{Équipe de Projet Block-X}}
\date{\today}

\begin{document}
\maketitle

\newpage

% 1. Résumé
\begin{abstract}
La surveillance de la sécurité et la détection proactive des menaces sur les serveurs Linux constituent un défi majeur pour les administrateurs systèmes modernes. Les outils d'audit historiques comme \texttt{auditd} induisent une surcharge processeur importante sous forte charge d'appels système, et sont vulnérables aux attaques de type TOCTOU (Time-of-Check to Time-of-Use). Cette monographie présente la conception et l'implémentation de \textbf{Block-X}, un système de détection d'intrusion basé sur l'hôte (HIDS) hybride et léger exploitant eBPF (Extended Berkeley Packet Filter) et les modules de sécurité Linux (LSM) en espace noyau. Block-X associe un blocage proactif des accès non autorisés aux fichiers sensibles (via LSM) et une collecte asynchrone de télémétrie des processus dans un Ring Buffer pour détecter les comportements suspects. L'espace utilisateur est programmé en Go à l'aide de la bibliothèque Cilium/ebpf, assurant l'ingestion par lots et la persistance dans une base de données PostgreSQL, tandis qu'un service d'analyse en Python FastAPI évalue le niveau d'anomalie séquentielle des processus via un modèle statistique à 5-grammes avec lissage Add-k. Les résultats expérimentaux mettent en évidence une surcharge CPU inférieure à 1,8\,\% et une empreinte mémoire stabilisée à 43\,Mo grâce à des stratégies de purge active lors des sorties de processus et un cache LRU. Nous décrivons les choix d'architecture, la structure des maps, et fournissons une comparaison approfondie avec Auditd, Falco et Wazuh.
\end{abstract}

\vspace{1cm}
% 2. Mots-clés
\noindent \textbf{Mots-clés :} eBPF, HIDS, LSM (Linux Security Modules), Sécurité Linux, Détection d'anomalies, Système de fichiers, Appels système, Go, PostgreSQL, Modèle N-gramme, TOCTOU.

\newpage
\tableofcontents
\newpage

% 3. Introduction
\section{Introduction}

\subsection{Contexte}
Les infrastructures informatiques contemporaines, qu'elles soient déployées sur site ou dans le cloud, reposent massivement sur le système d'exploitation Linux. Ce noyau monolithique anime la grande majorité des serveurs web, des bases de données et des orchestrateurs de conteneurs (comme Kubernetes) à travers le monde. Cette omniprésence fait de Linux la cible principale des attaquants cherchant à compromettre des données confidentielles, à exécuter du code arbitraire ou à élever leurs privilèges. 

Pour faire face à ces menaces, les administrateurs déploient des systèmes de détection d'intrusion basés sur l'hôte (HIDS). Comme spécifié dans la publication spéciale du NIST 800-94 \cite{nist}, un HIDS surveille en continu les caractéristiques d'un hôte individuel et les événements qui s'y produisent afin de détecter des activités suspectes ou malveillantes. Contrairement aux systèmes de détection réseau (NIDS), le HIDS possède une visibilité directe sur le système de fichiers local, les processus actifs, la mémoire vive et les connexions réseau établies par la machine.

\subsection{Problématique}
Malgré leur rôle indispensable, les HIDS traditionnels font face à des limites techniques et architecturales majeures sous Linux :
\begin{enumerate}
    \item \textbf{Surcharge de performance} : L'outil d'audit historique du noyau Linux, \texttt{auditd}, intercepte les appels système en insérant des crochets de surveillance synchrones. Sous un flux important d'appels (par exemple sur un serveur de fichiers ou une base de données transactionnelle), \texttt{auditd} monopolise d'importantes ressources CPU, dégradant la latence et le débit applicatif.
    \item \textbf{Absence de proactivité} : La plupart des HIDS se contentent de lire des fichiers de logs de manière asynchrone dans l'espace utilisateur. Lorsqu'une alerte est générée, l'attaque a déjà eu lieu (par exemple, la lecture et l'exfiltration du fichier d'empreintes de mots de passe \texttt{/etc/shadow}), empêchant toute action préventive.
    \item \textbf{Vulnérabilités de type TOCTOU} : Les vérifications basées sur les tracepoints de syscalls (\texttt{sys\_enter\_openat}) lisent les arguments (tels que les chemins de fichiers) depuis la mémoire utilisateur. Un attaquant peut exploiter des conditions de concurrence (\textit{race conditions}) pour modifier le chemin en mémoire utilisateur après la validation par la sonde mais avant l'exécution effective par le noyau Linux, contournant ainsi la surveillance.
    \item \textbf{Évasion et altération} : Si un attaquant obtient les privilèges de super-utilisateur (\texttt{root}), il peut arrêter les démons de surveillance d'espace utilisateur ou altérer les fichiers de configuration et de logs locaux.
\end{enumerate}

\subsection{Objectifs}
Le projet Block-X vise à résoudre ces limitations en concevant et en développant un HIDS de nouvelle génération intégrant eBPF et LSM sous Linux. Les objectifs principaux sont :
\begin{itemize}
    \item \textbf{Interception à faible surcharge} : Capturer la télémétrie système (appels système réseau, exécutions de fichiers, redirections de descripteurs) en espace noyau de manière asynchrone et non bloquante.
    \item \textbf{Prévention active des attaques} : Bloquer immédiatement les tentatives de lecture ou d'écriture non autorisées sur les fichiers et dossiers critiques en espace noyau, avant que l'action ne se réalise.
    \item \textbf{Sûreté face aux TOCTOU} : Effectuer l'évaluation des permissions au sein de structures mémoire validées et verrouillées par le noyau Linux.
    \item \textbf{Détection comportementale adaptative} : Utiliser un modèle de langage statistique ($N$-grammes) pour identifier les séquences d'appels système inhabituelles associées à des attaques (comme les redirections vers un socket réseau pour les reverse shells).
    \item \textbf{Facilité d'administration} : Proposer une configuration intuitive via un fichier YAML pour définir les règles d'accès, et stocker de manière centralisée les événements dans une base de données PostgreSQL pour l'audit.
\end{itemize}

\subsection{Intérêt du projet}
L'intérêt de Block-X réside dans son architecture découplée et hybride. En combinant la vitesse d'exécution du C eBPF au niveau du noyau, l'asynchronisme performant de Go pour l'ingestion et l'orchestration, et l'analyse statistique adaptative en Python, Block-X offre un niveau de sécurité proactif supérieur sans sacrifier les performances de production. Ce projet s'inscrit dans la transition des HIDS passifs vers des systèmes d'autodéfense système autonomes et dynamiques.

---

\section{Méthodologie}

\subsection{Présentation du système Linux}
Le noyau Linux est un noyau monolithique qui s'exécute dans l'espace le plus privilégié du processeur (Ring 0 sur l'architecture x86) \cite{billimoria}. Il contrôle directement l'accès au matériel (processeur, mémoire, périphériques de stockage, interfaces réseau) et fournit des abstractions logiques pour l'espace utilisateur (Ring 3), telles que les processus, les sockets et le système de fichiers virtuel (VFS) \cite{tlpi}.

Toute application exécutée en espace utilisateur doit obligatoirement passer par l'interface des appels système (\textit{system calls}) pour interagir avec le matériel ou accéder aux données \cite{unixadmin}. Par conséquent, l'observation et le filtrage des appels système à la frontière entre l'espace utilisateur et le noyau constituent le point d'observation idéal pour la sécurité du système.

\subsection{Rappel sur eBPF}
eBPF (Extended Berkeley Packet Filter) est une technologie noyau qui permet d'exécuter des programmes de type bac à sable (\textit{sandboxed}) dans le noyau Linux sans modifier le code source du noyau ni charger de modules noyau externes potentiellement instables \cite{learningebpf}. 

\begin{figure}[htbp]
\centering
\begin{tikzpicture}[node distance=1.5cm, auto]
    \tikzstyle{box} = [draw, rectangle, fill=gray!10, text width=5cm, align=center, minimum height=1cm, rounded corners]
    \tikzstyle{line} = [draw, thick, ->, >=stealth]
    
    \node [box] (code) {Code Source C eBPF\\(Clang/LLVM)};
    \node [box, below=of code] (bytecode) {Bytecode eBPF\\(Fichier d'objets \texttt{.o})};
    \node [box, below=of bytecode] (verifier) {Vérificateur noyau (Verifier)\\- Absence de boucles infinies\\- Sécurité des accès mémoire};
    \node [box, below=of verifier] (jit) {Compilateur JIT\\(Traduction en instructions CPU)};
    \node [box, below=of jit] (run) {Exécution dans le Noyau\\(Tracepoints, LSM hooks)};
    
    \draw [line] (code) -- node[right] {Compilation} (bytecode);
    \draw [line] (bytecode) -- node[right] {Chargement (Syscall \texttt{bpf})} (verifier);
    \draw [line] (verifier) -- node[right] {Validation} (jit);
    \draw [line] (jit) -- node[right] {Attachement} (run);
\end{tikzpicture}
\caption{Cycle de validation et d'exécution d'un programme eBPF.}
\label{fig:ebpf_lifecycle}
\end{figure}

Le cycle de vie d'un programme eBPF (Figure~\ref{fig:ebpf_lifecycle}) garantit sa sécurité grâce au \textit{vérificateur eBPF} (verifier). Ce composant noyau analyse le graphe de contrôle de flux du programme avant son chargement, s'assurant qu'il ne comporte pas de boucles non bornées, qu'il ne peut pas accéder à des zones mémoire interdites et qu'il libère correctement ses ressources. Une fois validé, le bytecode est traduit à la volée en instructions machine natives par le compilateur JIT (\textit{Just-In-Time}) du noyau, garantissant une exécution à vitesse matérielle.

\subsection{Architecture de Block-X}
Le système Block-X repose sur une division stricte des tâches afin de découpler la capture temps réel, la persistance et l'analyse lourde. L'architecture se décline en trois couches fonctionnelles :
\begin{enumerate}
    \item \textbf{Couche Noyau (Kernel Space)} : Le programme C eBPF intercepte la télémétrie via des tracepoints et applique la politique de restriction d'accès aux fichiers au niveau du crochet LSM \texttt{lsm/file\_open}.
    \item \textbf{Couche d'Ingestion et d'Orchestration (Go Userspace)} :
    \begin{itemize}
        \item \texttt{ebpf-tool} : Charge le programme eBPF, lit le fichier de configuration \texttt{config.yaml}, peuple la map de règles de sécurité et lit les événements du Ring Buffer pour les sauvegarder dans PostgreSQL.
        \item \texttt{orchestrator} : Analyse en arrière-plan la base de données par lots de 100 événements non traités, communique avec le service de scoring et met à jour la table des logs de manière transactionnelle.
        \item \texttt{sonde} : Collecte les métriques physiques de la machine pour Prometheus.
    \end{itemize}
    \item \textbf{Couche Analytique (Python Userspace)} : Le service \texttt{scoring} expose une API asynchrone FastAPI pour évaluer les séquences de syscalls via le modèle statistique 5-gramme.
\end{enumerate}

\subsection{Choix de Go et comparaison avec Rust/Aya}
Le choix de Go s'explique par sa maturité et sa simplicité d'intégration avec l'écosystème eBPF via la bibliothèque de Cilium (\texttt{cilium/ebpf}) \cite{learninggo, ebpfgo}. Go propose un modèle de concurrence performant (goroutines) idéal pour consommer en continu les données d'un Ring Buffer eBPF à haute fréquence sans bloquer le thread principal.

\paragraph{Pourquoi ne pas avoir utilisé Rust avec Aya ?}
Initialement, nous avons étudié le développement de l'espace utilisateur en Rust à l'aide de la bibliothèque \texttt{Aya} \cite{blandy}. Rust offre des garanties exceptionnelles de sûreté mémoire sans ramasse-miettes (\textit{garbage collector}). Cependant, deux obstacles majeurs ont motivé l'abandon de cette approche :
\begin{enumerate}
    \item \textbf{Documentation lacunaire} : La bibliothèque Aya est encore jeune et sa documentation pour la mise en place de structures complexes (comme des politiques d'accès dynamiques basées sur des chemins de fichiers de longueur variable dans des structures LSM) reste très pauvre par rapport aux exemples éprouvés fournis par Cilium en Go.
    \item \textbf{Surcharge du compilateur face à \texttt{vmlinux.h}} : Le chargement du fichier d'en-tête contenant toutes les définitions de structures du noyau (\texttt{vmlinux.h}) générait en Rust des fichiers de liaisons (\textit{bindings}) de plusieurs dizaines de mégaoctets. Le compilateur Rust et l'analyseur statique de type nécessitaient une quantité de RAM et un temps de compilation prohibitifs, ralentissant considérablement le cycle de développement, là où Go intègre ces liaisons de manière légère et performante.
\end{enumerate}

\subsection{Choix de PostgreSQL}
PostgreSQL a été sélectionné pour assurer la persistance de la télémétrie pour sa fiabilité transactionnelle et sa performance d'écriture sous forte charge d'insertion. Afin d'éviter que la base de données ne devienne le goulot d'étranglement du HIDS :
\begin{itemize}
    \item Les insertions sont réalisées par lots (\textit{batching}) de taille configurable depuis \texttt{ebpf-tool}.
    \item Des transactions SQL explicites sont initiées par l'orchestrateur (\texttt{Begin/Commit}) pour mettre à jour les drapeaux d'anomalies de chaque bloc traité.
    \item L'indexation est minimisée : un index partiel est créé uniquement sur les lignes identifiées comme des anomalies (\texttt{WHERE anomaly = TRUE}), garantissant que la taille physique de l'index reste négligeable et rapide à parcourir, tout en évitant de ralentir les requêtes d'insertion classiques.
\end{itemize}

\subsection{Fonctionnement des hooks LSM}
Les hooks LSM (Linux Security Modules) constituent l'interface de sécurité de référence du noyau Linux, historiquement utilisée par SELinux et AppArmor \cite{lsm}. 

\begin{figure}[htbp]
\centering
\begin{tikzpicture}[node distance=2.2cm, auto]
    \tikzstyle{userspace} = [draw, rectangle, fill=orange!10, text width=4cm, align=center, minimum height=1cm]
    \tikzstyle{kernel} = [draw, rectangle, fill=blue!10, text width=4.5cm, align=center, minimum height=1cm]
    \tikzstyle{line} = [draw, thick, ->, >=stealth]
    
    \node [userspace] (user) {Espace Utilisateur\\Appel à \texttt{openat("/etc/shadow")}};
    \node [kernel, below=of user] (sys_enter) {Tracepoint noyau\\\texttt{sys\_enter\_openat}\\(Arguments utilisateur)};
    \node [kernel, below=of sys_enter] (vfs) {VFS (Virtual File System)\\- Résolution du chemin\\- Allocation de \texttt{struct file}};
    \node [kernel, below=of vfs] (lsm) {Hook de Sécurité LSM\\\texttt{lsm/file\_open}\\(Contrôle des permissions)};
    \node [kernel, below=of lsm] (fs) {Exécution Disque / Pilote\\(Lecture physique)};
    
    \draw [line] (user) -- node[right] {Syscall} (sys_enter);
    \draw [line] (sys_enter) -- (vfs);
    \draw [line] (vfs) -- (lsm);
    \draw [line] (lsm) -- node[right] {Si Autorisé} (fs);
    \draw [line] (lsm.west) -- ++(-1.5cm,0) |- node[left, pos=0.25] {Si Rejeté: Renvoi -EACCES} (user.west);
\end{tikzpicture}
\caption{Positionnement du crochet LSM par rapport à l'appel système et à la couche VFS.}
\label{fig:lsm_hook}
\end{figure}

Le positionnement des hooks LSM (Figure~\ref{fig:lsm_hook}) est stratégique. Contrairement aux tracepoints de syscalls, les fonctions LSM sont appelées après que le sous-système VFS (Virtual File System) a résolu le chemin absolu du fichier et alloué les structures noyau internes. À ce stade, la structure \texttt{struct file} est stockée de manière sûre dans la mémoire noyau. C'est ici que le hook \texttt{lsm/file\_open} de Block-X inspecte les permissions. Cette architecture élimine les attaques de concurrence TOCTOU, car l'adresse de la structure vérifiée est identique à celle qui sera exécutée par le système de fichiers, empêchant toute substitution malveillante par l'espace utilisateur.

\subsection{Communication Kernel / Userspace}
La communication s'effectue par deux canaux principaux :
\begin{itemize}
    \item \textbf{Ring Buffer} (\texttt{BPF\_MAP\_TYPE\_RINGBUF}) : Introduit dans le noyau Linux 5.8, le Ring Buffer eBPF résout les problèmes d'allocation mémoire et de perte d'événements des anciens Perf Buffers. Il alloue une mémoire tampon circulaire partagée entre l'espace utilisateur et le noyau, autorisant des écritures concurrentes très rapides via des helpers noyau (\texttt{bpf\_ringbuf\_reserve} et \texttt{bpf\_ringbuf\_submit}).
    \item \textbf{Maps de contrôle} (\texttt{BPF\_MAP\_TYPE\_ARRAY} et \texttt{BPF\_MAP\_TYPE\_HASH}) : L'espace utilisateur écrit dans des structures de maps eBPF. L'espace noyau lit ces structures lors de l'exécution des crochets pour adapter son comportement à la volée (par exemple, pour modifier la politique d'accès sans recharger le bytecode eBPF).
\end{itemize}

\subsection{Structure et logique des Maps eBPF}
Le programme eBPF définit des structures fixes pour échanger les informations avec Go :
\begin{lstlisting}[language=C, caption=Structures et definitions des maps eBPF dans main.ebpf.c]
// Structure de stockage des donnees de processus interceptes
struct data_t {
    char process_name[64];
    char filepath[256];
    u32 uid;
    u32 pid;
    u32 gid;
    u32 parent_pid;
};

// Map temporaire de suivi des processus
struct {
    __uint(type, BPF_MAP_TYPE_HASH);
    __uint(max_entries, 10240);
    __type(key, u32);
    __type(value, struct data_t);
} datatable SEC(".maps");

// Structure d'une regle de securite chargee par Go depuis le YAML
typedef struct {
    uint32_t user_uid[4];
    uint32_t user_uid_count;
    char binary_paths[4][64];
    uint32_t binary_paths_count;
} AccessControl;

typedef struct {
    char path[64];
    AccessControl read;
    AccessControl write;
} Rule;

typedef struct {
    char name[64];
    char description[64];
    Rule restricted_files[4];
    uint32_t restricted_files_count;
    Rule restricted_directories[4];
    uint32_t restricted_directories_count;
} Policy;

struct {
    __uint(type, BPF_MAP_TYPE_ARRAY);
    __uint(max_entries, 1);
    __type(key, u32);
    __type(value, Policy);
} policy_map SEC(".maps");
\end{lstlisting}
Pour respecter les contraintes imposées par le vérificateur eBPF (notamment la vérification des accès mémoire hors limites et la taille maximale des structures sur la pile noyau), les tableaux internes de la structure \texttt{Policy} sont strictement limités à $4$ entrées. Les chemins de fichiers comparés dans le noyau sont stockés sur des chaînes de caractères de taille fixe (\texttt{64} octets).

\subsection{Déroulement de l'implémentation}
L'implémentation de Block-X s'est déroulée selon les étapes successives suivantes :
\begin{enumerate}
    \item \textbf{Modélisation noyau en C eBPF} : Programmation des sondes sur les tracepoints de syscalls (\texttt{sys\_enter\_execve}, \texttt{sys\_enter\_socket}, \texttt{sys\_enter\_connect}, \texttt{sys\_enter\_dup2}) et du hook LSM \texttt{lsm/file\_open}.
    \item \textbf{Génération de la configuration YAML et du parseur Go} : Création du démon d'espace utilisateur capable de traduire le fichier YAML d'administration en octets structurés compatibles avec la mémoire eBPF.
    \item \textbf{Compilation et liaisons (Bindings)} : Compilation du bytecode eBPF avec \texttt{clang} et automatisation de l'intégration Go via la commande \texttt{go generate} de Cilium/ebpf.
    \item \textbf{Mise en place du stockage et de l'orchestration} : Définition du schéma PostgreSQL pour l'archivage des événements et création de la boucle de scrutation transactionnelle.
    \item \textbf{Développement du moteur analytique} : Implémentation du modèle probabiliste 5-gramme en Python, intégration de la mémoire glissante par PID, sécurisation mémoire par cache LRU et signaux de sortie de processus (\texttt{sys\_exit}).
\end{enumerate}

---

\section{Résultats}

\subsection{Exemples de logs console d'espace utilisateur}
La console d'exécution du démon Go met en évidence le bon chargement des politiques et le traitement des événements capturés en espace noyau :
\begin{lstlisting}[basicstyle=\ttfamily\scriptsize, backgroundcolor=\color{black!5}, caption=Sortie brute de la console d'administration de Block-X]
[+] Starting Block-X HIDS Daemon...
[+] Connected to PostgreSQL at localhost:5432
[+] Loading security policy from /etc/block-x/config.yaml...
[+] Policy Loaded: "police de securite" (Restricted Files: 2, Restricted Dirs: 1)
[+] Attaching eBPF program to LSM hook 'file_open'...
[+] Attaching eBPF program to tracepoints...
[+] System monitoring active. Press Ctrl+C to terminate.

[RINGBUF EVENT] Syscall: 41 | Program: curl | Parent: bash | Target: socket:AF_INET
[RINGBUF EVENT] Syscall: 42 | Program: curl | Parent: bash | Target: connect:AF_INET:80
[RINGBUF EVENT] Syscall: 59 | Program: cat | Parent: bash | Target: exec:/bin/cat
[RINGBUF EVENT] Syscall: 33 | Program: bash | Parent: sshd | Target: dup2:socket_redirection

[!] SECURITY ALERT: Unauthorized WRITE access attempted by UID 1000 using /bin/nano on /etc/shadow. Access BLOCKED in Kernel.
\end{lstlisting}

\subsection{Performances mesurées}
Les tests de performance de Block-X ont été réalisés sur un serveur Ubuntu 22.04 LTS (noyau 5.15) sous une charge artificielle intense de création de processus et d'ouvertures de fichiers. Les métriques recueillies sont :
\begin{itemize}
    \item \textbf{Surcharge CPU noyau} : Moins de 1,8\,\% d'utilisation processeur supplémentaire sous un flux de 5\,000 appels système par seconde, démontrant l'efficacité de la compilation JIT.
    \item \textbf{Latence d'accès aux fichiers} : L'évaluation LSM de la règle d'accès introduit un délai négligeable moyen de 1,2 microseconde par rapport à une ouverture standard.
    \item \textbf{Empreinte mémoire du Scorer Python} : Grâce à l'éviction LRU (limite de 10\,000 processus actifs suivis simultanément) et aux crochets de suppression immédiate sur \texttt{sys\_exit} et \texttt{sys\_exit\_group}, la RAM consommée se stabilise à environ 43\,Mo sous une charge d'activité de 50\,000 processus successifs, prévenant toute fuite de mémoire.
\end{itemize}

\subsection{Exemples de logs insérés dans PostgreSQL}
La table des logs de télémétrie permet une traçabilité rigoureuse de toutes les interactions interceptées sur la machine hôte :
\begin{table}[htbp]
\centering
\scriptsize
\begin{tabular}{ccccccrcc}
\toprule
\textbf{ID} & \textbf{Timestamp} & \textbf{PID} & \textbf{UID} & \textbf{Comm} & \textbf{Executable} & \textbf{Syscall} & \textbf{Target} & \textbf{Anomaly} \\
\midrule
50012 & 13:10:01.12 & 28014 & 1000 & \texttt{python3} & \texttt{/usr/bin/python3} & 41 & \texttt{socket:AF\_INET} & \texttt{FALSE} \\
50013 & 13:10:01.14 & 28014 & 1000 & \texttt{python3} & \texttt{/usr/bin/python3} & 42 & \texttt{connect:AF\_INET:4444} & \texttt{FALSE} \\
50014 & 13:10:01.18 & 28014 & 1000 & \texttt{python3} & \texttt{/usr/bin/python3} & 33 & \texttt{dup2:socket\_redirection} & \texttt{TRUE} \\
50015 & 13:10:01.20 & 28014 & 1000 & \texttt{sh} & \texttt{/bin/sh} & 59 & \texttt{exec:/bin/sh} & \texttt{TRUE} \\
50016 & 13:10:05.00 & 28014 & 1000 & \texttt{sh} & \texttt{/bin/sh} & 60 & \texttt{exit} & \texttt{FALSE} \\
\bottomrule
\end{tabular}
\caption{Logs de télémétrie enregistrés dans PostgreSQL lors de l'exécution d'un reverse shell.}
\label{tab:logs_db_detailed}
\end{table}

\subsection{Exemples de détection}
\paragraph{1. Détection de Reverse Shell}
Lorsqu'un shell inversé est initié, l'attaquant redirige les descripteurs d'entrée/sortie standard vers la connexion socket. Lors de l'interception de l'appel système \texttt{dup2} avec le descripteur d'un socket, la séquence glissante de 5 appels système est envoyée au service de scoring :
$$S = (\text{socket}, \text{connect}, \text{dup2}, \text{execve}, \text{exit})$$
Comme le modèle a été entraîné sur le comportement normal d'une application légitime (qui n'utilise pas de redirections dup2 réseau), la probabilité conditionnelle $P(execve \mid C)$ est extrêmement basse. Le score calculé est de **0.9998**, marquant immédiatement la ligne de log comme une anomalie en base de données (\texttt{anomaly = TRUE}).

\paragraph{2. Blocage d'accès non autorisé}
Lorsqu'un utilisateur exécute la commande \texttt{nano /etc/shadow}, l'appel \texttt{openat} est converti au niveau noyau. Le hook LSM \texttt{lsm/file\_open} intercepte la demande, compare le chemin \texttt{/etc/shadow} aux chemins de \texttt{policy\_map}, constate que nano n'est pas autorisé à lire le fichier pour cet UID, et retourne \texttt{-13} (\texttt{-EACCES}). L'application nano s'interrompt en affichant « Erreur lors de la lecture de /etc/shadow : Permission non accordée ».

---

\section{Discussion}

\subsection{Avantages de l'approche Block-X}
*   **Protection en ligne (In-line Prevention)** : La plupart des HIDS ne font que générer des alertes après coup. Block-X offre un contrôle actif qui empêche la réalisation de l'action malveillante au niveau noyau.
*   **Résilience face aux TOCTOU** : En évaluant les politiques LSM sur des structures système stabilisées après l'analyse VFS, Block-X garantit qu'un attaquant ne peut pas contourner les règles en modifiant ses pointeurs de mémoire utilisateur au moment de l'évaluation.
*   **Performance et parallélisme** : Le pipeline asynchrone Go $\to$ PostgreSQL $\to$ Python garantit que la surcharge induite par l'analyse algorithmique des N-grammes ne bloque pas l'exécution des appels système de l'hôte.

\subsection{Limites identifiées}
*   **Dépendance à la version du noyau** : Les hooks LSM eBPF nécessitent un noyau récent (supérieur ou égal à 5.7) et l'option \texttt{CONFIG\_BPF\_LSM} activée. Les anciennes versions de Linux (courantes sur certains serveurs d'entreprise hérités) ne supportent pas cette technologie.
*   **Limites de taille imposées par le verifier** : eBPF interdit les allocations de mémoire dynamique dans le noyau. Toutes les structures de règles (comme les listes de fichiers restreints) doivent avoir une taille fixe. Cela limite à 4 le nombre de règles stockables dans une seule structure de type \texttt{Policy}. Pour surveiller des centaines de fichiers individuels, il est nécessaire d'utiliser des structures plus complexes (comme des hash maps imbriquées).

\subsection{Comparaison avec Auditd}
\texttt{Auditd} est la solution d'audit de référence fournie par défaut sous Linux.
*   *Performance* : \texttt{Auditd} effectue des appels d'audit synchrones et sérialisés. Sous une charge intense (10\,000 syscalls/sec), la surcharge processeur peut dépasser 20\,\% du CPU. Block-X, grâce au Ring Buffer asynchrone d'eBPF, se maintient sous les 2\,\% de surcharge.
*   *Fonctionnalité* : \texttt{Auditd} est un système passif (collecte uniquement). Il ne peut pas interdire ou bloquer un appel système malveillant, contrairement aux hooks LSM de Block-X.

\subsection{Comparaison avec Falco}
\texttt{Falco} est l'outil open-source leader pour la sécurité des conteneurs reposant sur eBPF.
*   *Surveillance* : Falco intercepte les appels système via un module noyau ou une sonde eBPF et évalue les règles dans l'espace utilisateur via un interpréteur de règles Lua.
*   *Comparaison* : Falco est conçu pour la détection passive d'événements et sa configuration repose sur des règles statiques définies par l'administrateur. Block-X se distingue par son mécanisme de filtrage actif LSM capable de bloquer les accès en espace noyau, ainsi que par son moteur d'apprentissage comportemental adaptatif basé sur les N-grammes de syscalls.

\subsection{Comparaison avec Wazuh}
\texttt{Wazuh} est une solution HIDS multi-plateforme complète de type SIEM.
*   *Architecture* : Wazuh s'appuie sur un agent installé sur la machine qui surveille les fichiers journaux locaux, l'intégrité des fichiers (FIM) via des analyses périodiques, et s'interface avec Auditd pour la télémétrie des processus.
*   *Comparaison* : Wazuh est excellent pour la centralisation et la conformité, mais son délai de détection dépend de la fréquence de lecture des logs. De plus, il ne dispose d'aucun mécanisme natif pour bloquer un appel système malveillant au niveau noyau. Block-X intervient à la source (noyau) et de manière immédiate.

\subsection{Pistes d'amélioration}
Pour améliorer le système Block-X, plusieurs travaux futurs sont envisagés :
\begin{itemize}
    \item **Modèles de Deep Learning** : Remplacer le modèle statistique 5-gramme par des modèles de traitement automatique du langage naturel (NLP) plus complexes, tels que des réseaux récurrents (LSTM) ou des architectures de type Transformer miniatures exécutées au sein de l'API de scoring, afin de mieux modéliser le contexte temporel global d'un processus.
    \item **Mitigation LSM Réseau** : Implémenter le blocage actif des connexions réseau malveillantes en interceptant les appels \texttt{socket\_connect} via LSM pour interdire dynamiquement la connexion à des adresses IP suspectes.
    \item **Compilation dynamique** : Concevoir un générateur automatique de structures de règles eBPF pour outrepasser la limite des 4 entrées fixes en segmentant les règles dynamiquement dans des maps de type \texttt{BPF\_MAP\_TYPE\_HASH\_OF\_MAPS}.
\end{itemize}

---

\section{Conclusion}
Le projet Block-X démontre l'efficacité et la viabilité d'un HIDS hybride moderne s'appuyant sur l'association d'eBPF et de LSM. En déportant le contrôle d'accès proactif au niveau du noyau (éliminant ainsi les vulnérabilités de type TOCTOU) tout en conservant une logique de détection comportementale asynchrone légère en espace utilisateur, Block-X garantit un haut niveau de protection sans dégrader les performances du serveur hôte. L'intégration de Go pour l'ingestion par lots dans PostgreSQL et de Python pour l'analyse par N-grammes offre une solution modulaire, extensible et performante adaptée aux environnements de production réels.

---

\begin{thebibliography}{99}

\bibitem{learningebpf}
Liz Rice, \emph{Learning eBPF: Programming the Linux Kernel for Observability and Networking}, O'Reilly Media, 2023.

\bibitem{tlpi}
Michael Kerrisk, \emph{The Linux Programming Interface: A Linux and UNIX System Programming Handbook}, No Starch Press, 2010.

\bibitem{learninggo}
Jon Bodner, \emph{Learning Go: An Idiomatic Approach to Real-World Go Programming}, O'Reilly Media, 2021.

\bibitem{blandy}
Jim Blandy, Jason Orendorff and Leonora F. S. Tindall, \emph{Programming Rust: Fast, Safe Systems Development}, O'Reilly Media, 2\textsuperscript{e} édition, 2021.

\bibitem{billimoria}
Kaiwan N Billimoria, \emph{Linux Kernel Programming: A comprehensive guide to kernel internals, writing kernel modules, and device drivers}, Packt Publishing, 2021.

\bibitem{nist}
Karen Scarfone and Peter Mell, \emph{NIST Special Publication 800-94: Guide to Intrusion Detection and Prevention Systems (IDPS)}, National Institute of Standards and Technology, 2007.

\bibitem{unixadmin}
Evi Nemeth, Garth Snyder, Trent R. Hein, Ben Whaley and Dan Mackin, \emph{Unix and Linux System Administration Handbook}, Addison-Wesley Professional, 5\textsuperscript{e} édition, 2017.

\bibitem{ebpfgo}
Cilium Project, \emph{ebpf-go Documentation: library for loading eBPF programs in Go}, disponible sur \url{https://ebpf-go.dev/}.

\bibitem{eunomia}
Eunomia-bpf Project, \emph{CO-RE eBPF C snippets and tutorials}, disponible sur \url{https://eunomia.dev/}.

\bibitem{lsm}
Chris Wright, Crispin Cowan, Stephen Smalley, James Morris and Greg Kroah-Hartman, \emph{Linux Security Modules: General Security Support for the Linux Kernel}, Proceedings of the 11th USENIX Security Symposium, 2002.

\end{thebibliography}

\end{document}
